\chapter{Conclusion and Future Work}
\label{chap:conclusion}

% -----------------------------------------------------------------------
% TO WRITE: Summarize what you did, what you found, and where to go next.
% Write this chapter last, once all results are in.
% Suggested structure below — remove these comments as you write.
% -----------------------------------------------------------------------

\section{Summary of Contributions}
% Briefly restate the research question and summarize what this thesis did:
% - The comparative framework you built
% - The dataset and experimental setting
% - The four methods evaluated

\section{Key Findings}
% Summarize the main results:
% - Did adaptive methods outperform A/B testing? By how much?
% - Which method achieved the best trade-off?
% - What did segment-level analysis reveal about personalization?
% - Was RL viable given the available data?

\section{Practical Recommendations}
% Translate your findings into actionable advice for marketing teams:
% - When is Thompson Sampling the right choice?
% - When are contextual bandits worth the added complexity?
% - What data volume is needed before RL becomes viable?
% - What GDPR and transparency considerations apply?

\section{Limitations}
% Honestly discuss what this thesis could not show or prove:
% - Reliance on simulated rather than fully live deployment
% - Single dataset from one platform (ProvenExpert)
% - Limited variant space tested

\section{Future Work}
% Suggest concrete directions for follow-up research:
% - Live A/B vs. bandit comparison with real traffic splits
% - Incorporating generative AI (LLMs) for automatic variant generation
% - Extending RL to longer multi-step user journeys
% - Applying the framework to other digital channels (push, SMS, in-app)
